% Chapter 2, Figure 51: aerodynamic techniques for inducing roll (scan: figures/ch2/fig51.png).
% One rocket, nose up, drawn three times with a different fin detail; coordinates are tens of the scan's
% pixels (150 per inch; tail at y = 0, axis at x = 0 of each), drawn 1.65 times the printed size.
% Rocket as measured on the three copies: length 21.5, nose 5.45, body diameter 1.85; fins root chord 3.2,
% tip chord 1.25, span 2.15, trailing edge square at the tail.
%   Airfoiled fins: the fin facing the reader, edge-on, shows its root section: a flat-bottomed airfoil (the
%     NACA four-digit thickness form, 12 percent thick, flat side to the right, as printed) along the root chord.
%   Canted fins: the edge-on fin slants 13.4 degrees across the axis (0.76 over the root chord, as printed);
%     the side fins' roots are the two slanted lines on the body (0.92 to 0.6 of the radius, as printed).
%   Spinnerons: a tab on the trailing edge of each fin, chord 0.85, bent 33 degrees: the side fins' tabs
%     foreshortened (0.71 deep, inner edge cut back, as printed); the edge-on fin's tab and that of the fin
%     behind it, bent opposite ways, show below the tail.
% The edge-on fins (each seen along its span, so showing its section) are the subject of the figure and are
% drawn in the outline weight, as heavy as the body (the 1973 art inks them heavier still). The three
% labels share one baseline.
% Overlay on the scan (each rocket aligned separately, redraw ink against the scan's ink): 100% of the
% redraw's ink within 3 px of the scan's for all three rockets (mean 0.1-0.2 px).
\documentclass[11pt]{standalone}
\usepackage{tamrfig}
\begin{document}
\begin{tikzpicture}[x=0.28cm, y=0.28cm,
    declare function={yt(\s)=5*0.12*(0.2969*sqrt(\s)-0.1260*\s-0.3516*\s^2+0.2843*\s^3-0.1036*\s^4);}]
  \tikzset{edge fin/.style={outline, fill=white}}
  \def\L{21.5}\def\R{0.925}\def\cr{3.2}
  % \model{<x of the axis>}{<label>}: the rocket, tail at y = 0, and its label below
  \newcommand{\model}[2]{%
    \pic[rotate=-90] at (#1,\L) {rocket={length=\L, diameter=1.85, nose length=5.45, root chord=\cr,
        tip chord=1.25, span=2.15, sweep=1.95, centerline=false}};
    \node[font=\small, anchor=base, inner sep=0pt] at (#1,-2.75) {#2};}

  % ---- airfoiled fins
  \model{0}{Airfoiled fins}
  \draw[edge fin] plot[domain=0:1, samples=50, smooth, variable=\t] ({0.15-2*\cr*yt(\t*\t)}, {\cr-\cr*\t*\t})
    -- cycle;

  % ---- canted fins
  \model{16}{Canted fins}
  \begin{scope}[shift={(16,0)}]
    \draw[edge fin] (-0.38-0.07,\cr) -- (-0.38+0.07,\cr) -- (0.38+0.07,0) -- (0.38-0.07,0) -- cycle;
    \draw[outline] plot[domain=0:1, samples=12, smooth] ({-\R*sin(67-30*\x)}, {\cr*(1-\x)});
    \draw[outline] plot[domain=0:1, samples=12, smooth] ({\R*sin(37+30*\x)}, {\cr*(1-\x)});
  \end{scope}

  % ---- spinnerons
  \model{32}{Spinnerons}
  \begin{scope}[shift={(32,0)}]
    \draw[edge fin] (-0.08,0) rectangle (0.08,\cr);
    \foreach \sg in {1,-1}
      \draw[outline, xscale=\sg] (1.15,0) -- (1.55,-0.71) -- (3.075,-0.71) -- (3.075,0);
    \foreach \a in {33,-33}
      \draw[edge fin, rotate=\a] (-0.08,0) rectangle (0.08,-0.85);
  \end{scope}
\end{tikzpicture}
\end{document}
